\section{Discussion}

\subsection{Why the Compounding Hypothesis Did Not Hold}
The working hypothesis going into this study was that later edits would show progressively more collateral damage, because each one operates on an already-modified image. The data shows the opposite, at the one point where a significant difference exists at all. The most plausible explanation, consistent with the baseball-glove example above, is that a chain which collapses catastrophically early has, by construction of a CLIP-similarity metric, very little room left to register as further changed in subsequent steps. An image that is already almost entirely black cannot become much more different from its original self than it already is. This is a measurement ceiling effect masking a front-loaded failure, not evidence that later edits are genuinely gentler. Separating ``drift is front-loaded and severe'' from ``drift does not compound'' would need a metric sensitive to further degradation of an already-degraded image, which is a natural direction for follow-up work rather than something the CLIP-similarity approach used here can resolve as built.

\subsection{Why Region-Locking Outperformed Masked Conditioning}
The prediction going in was that masked conditioning, a stricter constraint applied at generation time, would outperform region-locking's after-the-fact correction, at some cost to edit quality at the region boundary. The opposite happened on the drift metric itself, and the likely reason is a property of how the metric and the mitigation interact rather than either mitigation being fundamentally weaker. Masked conditioning edits a padded crop, the target box plus a 32-pixel margin, deliberately included so the model has enough surrounding context to produce a coherent edit. The Drift Score, however, only excludes the raw target box from scoring, so masked conditioning's own padding ring is counted as drift even when nothing objectionable happened there. Region-locking's hard revert has no equivalent margin. This should not be read as proof that region-locking produces better-looking edits in general; the qualitative example in Fig.~\ref{fig:qualitative} shows region-locking has its own artifact, a visible, unblended rectangular seam near the glove, consistent with the boundary-cost intuition behind the original hypothesis. It simply does not show up as a drift-score cost the way that hypothesis predicted.

\subsection{A Shared Weak Point: Target-Region Identification}
Both mitigations, and the Drift Score itself, depend on CLIP correctly identifying which region an instruction targets. This is reliable for ``change X'' and ``remove X'' instructions, where the target already exists in the pre-edit image. It is not reliable for ``add X'' instructions, where nothing pre-existing can match well. Tracing one concrete failure directly: for the instruction ``add a cap on the boy's head,'' the final step of the worst-case chain, the identified target box covered 99.6\% of the frame. Masked conditioning's spatial constraint therefore gave no real protection for that step, since it was cropping to almost the entire image regardless. Region-locking faced the identical oversized box on the same step and produced a fine result anyway, most likely because diffusion sampling is stochastic with no fixed seed, so the two mitigations' generation calls for that step were independent random draws from the same model, and one landed conservative while the other did not. This is a reproducible limitation of the target-identification approach, not an implementation bug in either mitigation, and it affects roughly half of the object-level chains, those whose instructions use an ``add'' verb.

To check whether this single traced case reflects a dataset-wide pattern, every step across all 480 was split into ``add'' instructions and all others, and compared with an independent-samples Mann-Whitney U test, since add and non-add steps are not paired observations within the same chain (Table~\ref{tab:addverb}). Baseline and region-locking both show lower mean drift on add steps than on other steps, a difference that is only borderline significant ($p=0.052$ and $p=0.072$ respectively) rather than conclusive. Masked conditioning is the outlier: its add-step mean is numerically \emph{higher} than its other-step mean, the only condition of the three where this reverses, though the difference itself does not reach significance ($p=0.693$) at this sample size. The direction is consistent with the mechanism traced above, an oversized target box removing masked conditioning's spatial constraint specifically on add instructions, but the dataset-wide test should be read as corroborating, not confirming, the single traced example; a larger sample would be needed to establish the effect conclusively.

\begin{table}[H]
\caption{Per-Step Drift by Instruction Type}
\label{tab:addverb}
\centering
\renewcommand{\arraystretch}{1.05}
\setlength{\tabcolsep}{4pt}
\begin{tabular}{@{}lccc@{}}
\toprule
Condition & ``Add'' Mean & Other Mean & Mann-Whitney $p$ \\
\midrule
Baseline             & 0.068 & 0.078 & 0.052 \\
Region-Locking       & 0.006 & 0.009 & 0.072 \\
Masked Conditioning  & 0.026 & 0.022 & 0.693 \\
\bottomrule
\end{tabular}
\end{table}

\subsection{Limitations}
Six constraints bound the present study. The 60-image, 120-chain dataset is adequate for the paired tests reported here but should be read as an initial empirical study rather than a definitive benchmark. All results are specific to InstructPix2Pix; the Drift Score is designed to be model-agnostic, but that design goal was not tested empirically against a second editor within this project's scope. CLIP similarity is an imperfect proxy for human-perceived change; an interactive human-rating tool covering 15 stratified chains was built for this purpose but was not administered, an explicitly optional step skipped under the 30-day time budget, so the CLIP-based numbers here carry that caveat without an independent human check. Target-region identification through CLIP's top-1 similarity is a known weak point for ``add'' instructions, as shown above. Region boxes, not exact per-pixel segmentation masks, are used throughout for both scoring and mitigation, for simplicity and consistency, and this coarser spatial unit is the direct cause of masked conditioning's padding-margin effect. Finally, the attention-restricted editing mitigation considered during scoping was not implemented, since RQ3 has a clear answer without it.

\subsection{Reproducibility}
The full pipeline is organized as a small, dependency-light toolkit rather than a single notebook. Drift-score computation, region segmentation, and CLIP embedding each live in their own module with unit tests that exercise the scoring logic against synthetic vectors and mocked segmentation output, independent of any GPU or network access. The two mitigation strategies share the same target-identification code path as the scorer by construction, rather than by convention, which is what makes the mechanistic explanations in Sections~V-B and~V-C traceable to a specific function rather than to inferred behavior. Every image, edit chain, intermediate result, and statistical test output is version-controlled alongside the code that produced it, so any reported number in this paper can be regenerated from the committed dataset without re-running the GPU stages. The chain-construction and scoring approach is not specific to InstructPix2Pix; substituting a different editor requires only replacing the module that calls it.

\subsection{Future Work}
The most direct extension is evaluating the Drift Score against a second, more recent editing model, to test the model-agnosticism claimed but not empirically exercised here. Running the built human-perceptual check would validate the CLIP-based drift numbers against human judgment on the same 15 chains. A drift metric less prone to the ceiling effect identified in Section~V-A, one that can register further degradation of an already-degraded image, would let a future study separate gradual compounding from sudden, front-loaded failure with more confidence than the present metric allows.
